\section{Regular cosets on the actual curve branches}
\label{sec:sr-curve-regular-cosets}
This section proves the regular-element density needed for the
rank-three Taylor--Wiles construction.  Its hypotheses are the
local separation and residual character conditions of the selected
curve points.  Cyclotomic irreducibility by itself is not used as a
substitute for them.

Write \(A=k[[T]]\), \(F=k((T))\), and let \(\overline F\) be an
algebraic closure.  Here \(k\) is finite of characteristic
\(\ell\ge5\), as in the generic principal-character branch.
Let \(G\) be the compact profinite global group and let
\(G'\) be its closed cyclotomic kernel.  A curve representation has
entries in \(A\), continuous modulo every \(T^j\).
The local ordinary characters are denoted by
\(\psi_{v,i}:D_v\to A^\times\).
The residual global determinant is the product of the characters
\(\chi_i:G\to k^\times\).  All residual characters in this section
are characters of the full group \(G\), not merely of \(G'\).

\subsection{Local separation on every open subgroup}
\begin{srlemma}[No torsion in the principal units]
\label{lem:sr-curve-principal-units}
The group \(1+Tk[[T]]\) is torsion free.
Suppose a compact local subgroup \(N_v\subset D_v\) maps into \(G'\),
the ratios \(\psi_{v,i}/\psi_{v,j}\) have residue one on \(N_v\),
and no such ratio is identically one there.
For every open subgroup \(V\subset G\), there is a
\(g\in N_v\) mapping into \(V\cap G'\) such that
\[
       \psi_{v,i}(g)^m\ne\psi_{v,j}(g)^m
       \qquad(i\ne j,\ m\ge1).
\]
\end{srlemma}
\begin{proof}
If \(z\in Tk[[T]]\) is nonzero and
\(m=\ell^r s\), with \(\ell\nmid s\), then
\[
 (1+z)^m=(1+z^{\ell^r})^s
       =1+s z^{\ell^r}+\text{terms of higher \(T\)-order}.
\]
The displayed leading term is nonzero.  This proves the first claim.

Put \(W=N_v\cap \iota_v^{-1}(V)\), where
\(\iota_v:D_v\to G\) is the local map.
It is open and of finite index in \(N_v\).
A character ratio nontrivial on \(N_v\) stays nontrivial on \(W\):
a positive power of any chosen element of \(N_v\) lies in \(W\),
and the principal units have no torsion.
Its image on \(W\) is therefore infinite.
The kernel of that ratio is closed and has infinite index in \(W\).
For completeness, such a subgroup has Haar measure zero:
for every positive integer \(M\) there are \(M\) distinct cosets,
all of its measure, and their disjoint union has measure at most one.
There are only finitely many pairs \(i\ne j\).
Their kernels cannot cover \(W\), because a finite union of
measure-zero sets has measure zero.
Choose \(g\) outside this union.
Torsion freeness now proves every asserted inequality.
\end{proof}

At a selected nice curve point, these hypotheses hold with
\(N_v\) the prescribed local subgroup on which both the
residual character and the cyclotomic character are trivial.
The equal residual values follow from that definition.
The condition usually denoted \(\mathrm{Nice.sep}\) says precisely
that the differences of the two universal local characters do
not vanish on this subgroup after specialization.
A scalar twisting character cancels from every ratio.
Thus the preceding lemma applies to the actual twisted point,
on every open subgroup, before any regular detector is selected.
This deduction still requires a selected curve with that
nonvanishing property; it does not prove existence of such curves.

\subsection{The prime-rank restriction argument}
We give the elementary irreducibility argument that strengthens
the global type condition on these curves.

\begin{srlemma}[Characters and reduction]
\label{lem:sr-curve-character-matching}
Let \(r:G\to\mathrm{GL}_d(A)\), \(d<\ell\), have residual
characteristic polynomial
\(\prod_{i=1}^d(X-\chi_i(g))\) at every \(g\).
Suppose on a subgroup \(P\) its extension to \(\overline F\)
is a sum of characters \(\lambda_1,\ldots,\lambda_d\).
There is a single permutation of the indexed multisets such that
the reductions of the \(\lambda_i\) equal
\(\chi_i|_P\), after that permutation.
The same assertion holds over \(\overline F\) for two sums of
characters with equal characteristic polynomials.
\end{srlemma}
\begin{proof}
Extend the \(T\)-adic valuation to \(\overline F\).
Every \(\lambda_i(g)\) is integral because it is a root of a
monic characteristic polynomial over \(A\).
Applying the same argument to \(g^{-1}\) shows its inverse is
integral.  It is a unit of this valuation ring and has a
nonzero residue in an algebraic closure of \(k\).
Reduction therefore gives genuine characters of \(P\).

We recall the independence argument for distinct characters.
If a nonzero relation \(\sum_i a_i\theta_i(g)=0\) holds for
every \(g\), choose one with the fewest nonzero coefficients.
For \(\theta_1\ne\theta_j\), choose \(h\) with
\(\theta_1(h)\ne\theta_j(h)\).
Subtract \(\theta_j(h)\) times the relation at \(g\) from the
relation at \(hg\).  This gives a shorter nonzero relation,
a contradiction.  Thus distinct characters are linearly
independent.

Take traces of the two characteristic-polynomial identities.
Group together equal characters and apply this independence.
Their multiplicities are integers between zero and \(d<\ell\),
so equality in characteristic \(\ell\) is equality of those
integers.  The two indexed multisets therefore agree.
Choosing a bijection of equal entries gives one fixed
permutation, valid on the whole subgroup.
\end{proof}

\begin{srproposition}[Irreducibility on the required restrictions]
\label{prop:sr-curve-open-irreducibility}
Let \(d=2\) or \(3\), and suppose \(r:G\to\mathrm{GL}_d(A)\)
is absolutely irreducible over \(\overline F\).
Assume that, for one local group \(D\), its semisimplification
is the sum of \(d\) ordinary characters \(\psi_i\).
Assume also:
\begin{enumerate}
\item for every open subgroup \(V\subset G\), some element of
\(D\) maps into \(V\cap G'\) with all \(\psi_i\)-values distinct;
\item the residual determinant is a sum of global characters
\(\chi_i\), and two of their restrictions to \(D\) are unequal.
\end{enumerate}
Then \(r|_{V\cap G'}\) is absolutely irreducible for every
open subgroup \(V\subset G\).
\end{srproposition}
\begin{proof}
Replace \(V\) by its open normal core and put \(H=V\cap G'\).
It suffices to prove irreducibility on \(H\).

Here is the prime-dimension Clifford argument in the needed form.
Choose a minimal nonzero \(H\)-invariant subspace \(W\).
Every translate \(r(g)W\) is simple for \(H\), because \(H\)
is normal.  The sum of all these translates is a nonzero
\(G\)-invariant subspace, hence the whole representation.
Only finitely many translates are needed to span it.
A sum of simple modules is semisimple: adjoining a simple
module either has zero intersection with the sum already
formed or is contained in it.
Grouping the resulting simple modules into their isotypic
parts, \(G\) acts transitively on these parts.
Their simple dimension, their common multiplicity and
their number have product \(d\).
Since \(d\) is prime, either the restriction is irreducible,
or \(H\) acts by one scalar character, or it has exactly
\(d\) distinct one-dimensional character spaces, permuted
transitively by \(G\).
In the homogeneous one-dimensional case \(H\) acts by scalars
on the entire space; this also follows directly from
semisimplicity.

The scalar possibility contradicts the element in (1):
its eigenvalues are the distinct values of the ordinary
characters.  Suppose then that there are \(d\) distinct
lines, and write
\[
        \pi:G\to S_d,\qquad P=\ker\pi .
\]
The action of \(P\) on these lines is by characters
\(\lambda_1,\ldots,\lambda_d\).
On \(D\cap H\) the line characters are the restrictions of
the \(\psi_i\), in one fixed order, by
Lemma~\ref{lem:sr-curve-character-matching}.
They are distinct there, again by (1).
For \(t\in D\), each character \(\psi_i\) is unchanged by
conjugation by \(t\).
Conjugation therefore cannot permute the distinct characters
on \(D\cap H\).  It fixes every one of the \(d\) lines.
Thus \(D\subset P\).

Reduce the \(\lambda_i\) using
Lemma~\ref{lem:sr-curve-character-matching}.
They become the restrictions of the global \(\chi_i\), in
one fixed order.  A global character is invariant under
conjugation by every element of \(G\).
But the permutation action of \(G\) is transitive on the
lines, so conjugation transports any \(\lambda_i\) to any
\(\lambda_j\).
Their reductions are consequently all the same character
on \(P\).
Since \(D\subset P\), all the \(\chi_i|_D\) are equal,
contradicting (2).
Only the irreducible case remains.
\end{proof}

For the principal rank-three stable branch, the local residual
characters are a common scalar times
\(\overline\Psi,1,\overline\Psi^{-1}\).
The host condition supplies a local element with
\(\overline\Psi^2\ne1\), so the residual distinction in (2)
holds.  Lemma~\ref{lem:sr-curve-principal-units} supplies (1).
The same argument applies to the rank-two block on the
two-plus-one branch.  In the actual integral coordinates,
this block uses the outer indices \(0,2\), and its residual
characteristic polynomial is
\[
        (X-\chi_0(g))(X-\chi_2(g)).
\]
Indeed one of the two opposite off-diagonal entries is zero
after reduction, so their product in the two-by-two
determinant vanishes.  The residual distinction is now
\(\overline\Psi\ne\overline\Psi^{-1}\).
The middle line is a direct summand of the ordinary local
representation; its character is one of the three ordinary
Jordan--Hölder characters.  The two-dimensional summand has
the other two.  Their restrictions to every \(V\cap G'\)
are distinct by the preceding local separation.
Proposition~\ref{prop:sr-curve-open-irreducibility} therefore
makes the outer block absolutely irreducible on every such
restriction.  No irreducibility is asserted for the full
two-plus-one representation.

\subsection{Closed zero sets and principal minors}
\begin{srlemma}[Closedness for algebraic-coefficient linear forms]
\label{lem:sr-curve-linear-closedness}
For every \(\overline F\)-linear functional \(L\) on
\(M_d(\overline F)\), the set
\(\{g:L(r(g))=0\}\) is closed in \(G\).
\end{srlemma}
\begin{proof}
The finitely many coefficients of \(L\) lie in a
finite-dimensional \(F\)-subspace of \(\overline F\).
Choose an \(F\)-basis \(b_1,\ldots,b_s\) for their span and
write
\[
                  L(r(g))=\sum_{\nu=1}^s b_\nu L_\nu(r(g)),
\]
where the \(L_\nu\) have coefficients in \(F\).
The maps on the right take values in \(F\) and are continuous
for its \(T\)-adic topology: each is a finite linear
combination of the continuous entries in \(A\).
Linear independence of the \(b_\nu\) says \(L(r(g))=0\)
exactly when all \(L_\nu(r(g))=0\).
This is a finite intersection of closed sets.
\end{proof}

\begin{srlemma}[A matrix in the coset with nonzero minors]
\label{lem:sr-curve-coset-minors}
Let \(Q\) be a compact profinite group with a representation
\(s:Q\to\mathrm{GL}_d(\overline F)\).
Assume its restriction to every open subgroup is irreducible,
and all linear-functional zero sets in the preceding lemma
are closed.
For every \(y\in\mathrm{GL}_d(\overline F)\), there exists
\(b\in Q\) such that every diagonal entry of \(ys(b)\)
and of its inverse is nonzero.
If \(d\le3\), every principal minor of \(ys(b)\) is nonzero.
\end{srlemma}
\begin{proof}
For a nonzero vector \(v\) and a nonzero linear form \(f\),
the zero set of \(f(s(g)v)\) has empty interior.
Otherwise it contains a coset \(hW\), with \(W\) open in \(Q\).
The span of \(s(W)v\) is a nonzero \(W\)-invariant subspace,
hence the whole space.  The relation
\(f(s(h)s(w)v)=0\) for every \(w\in W\) would make
the nonzero form \(f\circ s(h)\) zero on the whole space.
This is impossible.
The zero set is closed by the linear-functional hypothesis.
Its complement is therefore open and dense.

Apply this to the \(d\) diagonal coefficients of \(ys(g)\).
For the inverse coefficients use \(g\mapsto g^{-1}\), a
homeomorphism of \(Q\), and the diagonal coefficients of
\(s(g)y^{-1}\).
There are only \(2d\) closed zero sets.
The intersection of their open dense complements is nonempty:
intersect the complements one at a time, keeping a nonempty
open set at each step.
A point \(b\) in that intersection has the required property.

For a one-by-one principal minor this is the diagonal assertion.
For \(d=3\), a two-by-two principal minor complementary to
index \(i\) equals
\(\det(ys(b))(ys(b)^{-1})_{ii}\), by the cofactor formula.
It is nonzero.  The full determinant is nonzero because
\(ys(b)\) is invertible.  The empty minor equals one.
The cases \(d\le2\) are included in the same argument.
\end{proof}

The hypotheses here hold after every change of basis over
\(\overline F\), since a matrix coefficient in the new basis is
still a linear functional in the old entries.
They also hold on \(V\cap G'\): every open subgroup of this
closed subgroup contains a subgroup of the form
\(V_1\cap G'\), with \(V_1\) open in \(G\).
This follows by choosing an open ambient identity neighborhood
inside the given relative open neighborhood and then a
profinite open subgroup inside it.
The preceding restriction proposition gives irreducibility.

\subsection{The quadratic and cubic discriminants}
\begin{srlemma}[Separating a small matrix by diagonal powers]
\label{lem:sr-curve-discriminant}
Let \(d\le3\), \(M\in\mathrm{GL}_d(K)\), for a field \(K\),
have all principal minors nonzero.
Let \(\alpha_1,\ldots,\alpha_d\in K^\times\) satisfy
\((\alpha_i/\alpha_j)^m\ne1\) for \(i\ne j\) and \(m\ge1\).
Then the characteristic polynomial of
\(M\operatorname{diag}(\alpha_i)^m\) is separable for some
integer \(m\ge0\).
\end{srlemma}
\begin{proof}
Let \(B\) be the finitely generated multiplicative subgroup
generated by the \(\alpha_i\).
Its quotient by its finite torsion subgroup is a free
abelian group.  Each \(\alpha_i/\alpha_j\) is nonzero in
this quotient.  Choose a rational linear functional on
\(B\otimes_{\mathbb Z}\mathbb Q\) which is nonzero on every
one of these finitely many differences.
Such a functional exists because finitely many proper
hyperplanes cannot cover a vector space over the infinite
field \(\mathbb Q\). In a fixed basis, evaluate the finitely
many nonzero defining linear forms at \((1,t,t^2,\ldots)\).
Each becomes a nonzero rational polynomial in \(t\); an integer
outside their finitely many root sets gives the required functional.  It gives an additive weight
\(w:B\to\mathbb Q\) with all \(w(\alpha_i)\) distinct.

We shall use the following elementary fact.
An exponential sum \(\sum_{\beta}c_\beta\beta^m\), with
distinct nonzero bases \(\beta\), cannot vanish for every
\(m\ge0\) unless all \(c_\beta\) are zero.
For \(s\) bases, the equations for \(m=0,\ldots,s-1\)
have Vandermonde determinant
\(\prod_{i<j}(\beta_j-\beta_i)\ne0\).
In particular a term with uniquely largest \(w\)-weight
cannot cancel with the other bases.

For \(d=0,1\), separability is immediate.
For \(d=2\), relabel so \(w_1>w_2\), and put
\(d_i=M_{ii}\), \(D=\det M\).
The discriminant is the exponential expression
\[
 (d_1\alpha_1^m+d_2\alpha_2^m)^2
                      -4D(\alpha_1\alpha_2)^m .
\]
Its uniquely highest-weight term is
\(d_1^2(\alpha_1^2)^m\), with nonzero coefficient.
It is nonzero for some \(m\), proving the claim.

For \(d=3\), relabel so \(w_1>w_2>w_3\).
Write \(p_{ij}\) for the principal two-by-two minors.
The elementary symmetric coefficients of the characteristic
polynomial are
\[
 e_1=\sum_i d_i\alpha_i^m,\qquad
 e_2=\sum_{i<j}p_{ij}(\alpha_i\alpha_j)^m,\qquad
 e_3=D(\alpha_1\alpha_2\alpha_3)^m .
\]
The discriminant of \(X^3-e_1X^2+e_2X-e_3\) is
\[
 \Delta=e_1^2e_2^2-4e_2^3-4e_1^3e_3
                       -27e_3^2+18e_1e_2e_3 .
\]
The largest possible weights in these five terms, in order,
are
\[
 \begin{gathered}
 4w_1+2w_2,\quad 3w_1+3w_2,\quad
 4w_1+w_2+w_3,\\
 2w_1+2w_2+2w_3,\quad 3w_1+2w_2+w_3 .
 \end{gathered}
\]
The first is strictly largest.
Within its term the unique largest monomial is
\[
              d_1^2p_{12}^2(\alpha_1^4\alpha_2^2)^m .
\]
Its coefficient is nonzero.  The Vandermonde fact therefore
makes \(\Delta\ne0\) for some \(m\).
This reasoning is valid also when some of the integers
\(4,18,27\) become zero in \(K\); the leading coefficient
has coefficient one.
Nonzero discriminant is equivalent to separability for a
monic polynomial, completing the proof.
\end{proof}

\subsection{Density on both actual rank-three branches}
\begin{srtheorem}[The curve-point regularity condition]
\label{thm:sr-curve-N3}
Consider a selected nice curve point of the generic principal
rank-three construction.  On its stable branch assume the
global representation is absolutely irreducible.
On its two-plus-one branch assume the actual outer
rank-two block is absolutely irreducible and the middle
coordinate is a character.
Then, on either branch, every open subset of \(G\)
meeting \(G'\) contains a \(\sigma\in G'\) for which
\(\det(X-r(\sigma))\) is separable.
\end{srtheorem}
\begin{proof}
Let \(O\) be such an open subset, choose \(t\in O\cap G'\),
and choose an open subgroup \(V\subset G\) with \(tV\subset O\).
Local separation supplies \(u\in V\cap G'\) whose full
ordinary eigenvalue tuple \(\alpha\) has torsion-free
pairwise ratios.  It has distinct eigenvalues, so
\(r(u)\) is diagonalizable over \(\overline F\).
Work in that eigenbasis.

On the stable branch,
Proposition~\ref{prop:sr-curve-open-irreducibility} and
Lemma~\ref{lem:sr-curve-coset-minors}, applied to
\(Q=V\cap G'\) and \(y=r(t)\), give \(b\in Q\)
such that \(M=r(t)r(b)\) has every principal minor nonzero.
Lemma~\ref{lem:sr-curve-discriminant} now gives an \(m\)
for which \(Mr(u)^m\) has separable characteristic polynomial.
It is \(r(tb u^m)\), and \(tb u^m\in O\cap G'\).

On the two-plus-one branch diagonalize the rank-two block
of \(r(u)\); its remaining scalar is the middle-line value.
The rank-two restriction is irreducible on every open subgroup
of \(Q\), as proved above.
Choose \(b\in Q\) so that the two diagonal entries of
the rank-two block of \(r(t)r(b)\) are nonzero.
The scalar block is nonzero because it is a character.
Every principal minor of the resulting three-by-three
block matrix is nonzero: the mixed two-by-two minors are
that scalar times a diagonal entry, the remaining
two-by-two minor is the invertible rank-two determinant,
and the full determinant is their product.
Apply Lemma~\ref{lem:sr-curve-discriminant} to the full
three-by-three matrix and its full free-ratio tuple.
The same element \(tb u^m\) works.

Separable characteristic polynomials descend from
\(\overline F\) to \(F\), since extension of fields preserves
the nonvanishing of their discriminant.
This proves exactly the regular-density condition
usually denoted \(N3\), before taking finite lists of
simultaneously regular detectors.
\end{proof}

The arithmetic content used here is the actual nonvanishing
of the local character differences on the selected curve,
the residual local distinction, and the representation type
of that curve.  Selecting curves with these properties and
placing them in the required common support remain the
separate curve-selection and support arguments.
The elementary density proof introduces no additional
automorphy assumption or fixed positive-characteristic
valuation on \(\overline F\).

For correspondence with the frozen development, the relevant
paths under \srcid{CurveDVRSpecialization} are
\srcid{NiceSep}, \srcid{IrrNormal}, \srcid{CliffPrime},
\srcid{NiceIrr}, \srcid{NiceIrrE},
\srcid{LinClosedOfAdic}, \srcid{RegCoset},
\srcid{RegCosetE}, and
\srcid{NiceSpecializationTwistSatisfiesThirdCondition}.
The closed-zero-set argument above replaces the more general
abstract exponential-sum/full-matrix construction by a direct
argument on the actual profinite curve points.


\subsection{The remaining two local nice-point conditions}
\begin{srproposition}[The distinguished local invariant space]
\label{prop:sr-curve-N4}
Let \(r:G\to\mathrm{GL}_n(A)\) have the prescribed residual
determinant. Suppose a distinguished local image \(D_0\)
has the following two properties:
all residual \(\chi_i|_{D_0}\) agree, and some \(g\in D_0\)
has \(\overline\epsilon(g)\ne1\).
Then
\[
 \{Y\in M_n(F):
      \epsilon(g)r(g)Y=Yr(g)\text{ for every }g\in D_0\}=0.
\]
The assertion holds for every point of that same residual
determinant and in every rank.
\end{srproposition}
\begin{proof}
The zero-rank case is immediate.
Use just the distinguished \(g\).
Write \(C=r(g)\), \(e=\epsilon(g)\), and let
\(c\in k^\times\subset A^\times\) be its common residual
character value.
Cayley--Hamilton applied after reduction gives
\((C-cI)^n\in T M_n(A)\).
Put \(N=C-cI\) and \(V=eC-cI\).
Modulo \(T\),
\[
                  V=(\overline e-1)cI+\overline e\,N.
\]
This is invertible: a nonzero scalar plus a commuting
nilpotent matrix has inverse given by a finite geometric
series. Its determinant is a unit of \(A\), so \(V\) is
invertible over \(A\).

Clear a common denominator of \(Y\), so \(Y\in M_n(A)\).
The assumed equation is \(VY=YN\).
Induction gives \(V^nY=YN^n\).
Write \(N^n=TQ\), with \(Q\in M_n(A)\); then
\[
                         Y=T V^{-n}YQ.
\]
Iterating this equality shows every entry of \(Y\) belongs
to \(T^jA\) for every \(j\).
Their intersection is zero, so \(Y=0\).
Clearing denominators was injective, which proves the
assertion over \(F\).
\end{proof}

For the generic principal packet over
\(K=\mathbb Q(\sqrt{-3})\), this local condition can itself
be arranged concretely.
Here \(3\nmid\ell N\), \(\ell\) is odd, and
\(\psi(-1)=(-1)^a\).
Choose a prime \(p\), outside any prescribed finite set, with
\[
                       p\equiv-1\pmod{\ell N},
                       \qquad p\equiv1\pmod3.
\]
The Chinese remainder theorem gives a reduced residue
class, and Dirichlet's theorem supplies infinitely many
such primes.
It splits in \(K\).
Its residual arithmetic Frobenius has
\(\overline\epsilon(p)=-1\) and
\[
           \overline\epsilon(p)^a\overline\psi(p)
                     =(-1)^a\overline\psi(-1)=1.
\]
The three residual characters are therefore all equal
on its full decomposition group, after the common scalar
link twist; inertia is trivial in the residual characters
at this prime.
At the initial imaginary-quadratic packet, taking that local
group for \(D_0\) gives exactly the hypotheses of
Proposition~\ref{prop:sr-curve-N4}.  When the arithmetic link
passes to a further CM tower, it must retain these two
local properties; the initial congruence calculation alone
does not establish that transfer, since a residue-degree
power can change the cyclotomic value.
The local level introduced to obtain this distinguished
place must subsequently be removed by the two-link argument.
No bound for the least such prime is used here.
For other residual packets the two displayed local
properties must be checked separately.

Finally, the ordered ordinary condition \(N1\) follows from
the same local separation used above.
Choose its free-ratio element in the full cyclotomic kernel.
The \(\psi_i\) are pairwise unequal there.
If \(\psi_i=\psi_j\epsilon\) with \(i<j\), their values
would also be equal on this element, because
\(\epsilon=1\) there.  This is a contradiction.
Consequently the selected rank-three curves satisfy \(N1\)
and \(N3\).  Whenever the selected link supplies the
explicit distinguished-place condition \(\mathrm{V0ok}\),
they also satisfy \(N4\) by the full matrix argument above. Their stable and two-plus-one type conditions are
the ones retained throughout the preceding arguments.

The general \(N4\) calculation corresponds to
\srcid{WeakShimuraTheory/ProjectiveLineFlags/ContinuousLinearClosed.lean},
\srcid{P1.n4V0}.  Its hypothesis is the explicit residual
local condition \srcid{V0ok}; a global no-self-twist or
no-Speh condition is not substituted for it.
